%% ============================================================================
\chapter{Limitations and Future Work}
\label{ch:futurework}
%% ============================================================================

\section{Limitations}

\paragraph{A single filesystem and a single cluster.}
All primary measurements were taken on GPFS on the Iris cluster.
The many-small-file pathology is general to parallel filesystems, but its magnitude
is not: Lustre, BeeGFS and object stores differ in metadata architecture, and the
threshold at which sharding becomes worthwhile will differ accordingly.
The classification rules are expressed as ratios rather than absolute times
specifically so that they might transfer, but that transfer is asserted rather than
demonstrated.

\paragraph{Contention is uncontrolled.}
Read latency on a shared filesystem depends on concurrent jobs.
Repeats and reporting of variance mitigate this, but do not eliminate it, and a
measurement taken on a quiet cluster is not the measurement a user will experience
on a busy one.
An adversarial evaluation --- deliberately co-scheduling I/O-heavy jobs --- would
strengthen the claim considerably and was not undertaken here.

\paragraph{One corpus.}
The design generalises to any workload with a fixed epoch permutation over many
small files, but it has been evaluated on one dataset of approximately 1~GB.
A corpus one or two orders of magnitude larger --- where staging cannot be a
full copy under any policy --- is where the adaptivity should matter most and where
it remains untested outside the synthetic capacity constraint of
Section~\ref{sec:ablation-capacity}.

\paragraph{Single-node scope.}
\pads{} stages to node-local scratch for a single training process.
Multi-node distributed training raises questions the present design does not
address: whether shards should be partitioned across ranks or replicated, how
staging interacts with a distributed sampler, and whether one rank's staging
traffic degrades another's read latency.

\paragraph{The classifier is a rule set.}
Bottleneck classification uses a small number of hand-chosen thresholds.
This was a deliberate choice --- an inspectable rule set is auditable and its
failures are diagnosable, which a learned policy is not --- but it means the
thresholds are the designer's judgement.
Section~\ref{sec:ablation-thresholds} reports their sensitivity rather than removing
the concern.

\section{Future work}

\paragraph{Decode offload.}
The classifier can already detect the case in which \Tdecode{} dominates, but its
only response is to raise the worker count.
GPU-side decode~\cite{DALI} or a cache of decoded, pre-augmented tensors would give
it a substantive action in that regime.
The trade-off is memory footprint against decode cost, and the profiler is already
measuring the quantities needed to make that choice.

\paragraph{Cross-job coordination.}
On a shared cluster, several jobs frequently train on the same corpus.
Mohan et al.~\cite{Mohan21} show substantial gains from coordinating caching and
eliminating redundant preprocessing across concurrent jobs.
Because \pads{} already materialises shards on node-local scratch, extending it to
publish and reuse staged shards across co-scheduled jobs is a natural next step.

\paragraph{Online rather than one-shot adaptation.}
The current design profiles once, at the start of the run, and commits.
Since contention varies over the lifetime of a job, periodic re-profiling with
hysteresis --- re-evaluating the policy every $N$ epochs, and switching only on a
sustained change --- would address the non-stationarity that motivated adaptivity in
the first place.
This is arguably the most direct extension of the present work.

\paragraph{Revisiting the accuracy question.}
Two findings from the reproduction bear on segmentation quality rather than
throughput and were not pursued here.
The first is that per-image IoU is strongly monotonic in the fraction of the crop
occupied by the site (Table~\ref{tab:by-area}), with tiles below 2\% site area
scoring near zero; a size-aware evaluation, or a size-stratified report, would give
a more informative picture than a single aggregate.
The second is that batch-normalisation statistics degrade badly at the batch sizes
imposed by consumer accelerators, which suggests that a normalisation scheme less
sensitive to batch size~\cite{Wu18} would make this workload substantially more
reproducible outside a cluster.
Neither is a data-staging question, but both affect who can reproduce this work.

\paragraph{Beyond archaeology.}
The workload characteristics that motivate \pads{} --- many small paired files, a
fixed permutation, a corpus that may or may not fit node-local scratch --- are
common across remote sensing, medical imaging and microscopy.
Evaluating on a second domain would establish whether the contribution is about this
dataset or about a class of pipelines.
